% !TEX program = lualatex
% Transparencias de Fundamentos de las Matemáticas II
% Clase 05: Algoritmo de Euclides e identidad de Bézout

\documentclass[aspectratio=169,11pt]{beamer}

\usepackage{fontspec}
\usepackage[spanish,es-nodecimaldot]{babel}
\usepackage{amsmath,amssymb,mathtools}
\usepackage{booktabs}
\usepackage{csquotes}
\usepackage{xcolor}

\usetheme{Madrid}
\usecolortheme{default}
\setbeamertemplate{navigation symbols}{}
\setbeamertemplate{footline}[frame number]

\definecolor{FdMBlue}{RGB}{20,55,100}
\definecolor{FdMRed}{RGB}{130,30,30}
\definecolor{FdMGreen}{RGB}{25,100,60}

\setbeamercolor{structure}{fg=FdMBlue}
\setbeamercolor{title}{fg=white,bg=FdMBlue}
\setbeamercolor{frametitle}{fg=white,bg=FdMBlue}
\setbeamercolor{block title}{fg=white,bg=FdMBlue}
\setbeamercolor{block body}{bg=blue!3}
\setbeamercolor{alerted text}{fg=FdMRed}

\setmainfont{Latin Modern Roman}
\setsansfont{Latin Modern Sans}
\setmonofont{Latin Modern Mono}

\newcommand{\N}{\mathbb{N}}
\newcommand{\Z}{\mathbb{Z}}
\newcommand{\Q}{\mathbb{Q}}
\newcommand{\R}{\mathbb{R}}
\newcommand{\C}{\mathbb{C}}
\newcommand{\Pow}{\mathcal{P}}
\newcommand{\id}{\operatorname{id}}
\newcommand{\mcd}{\operatorname{mcd}}
\newcommand{\mcm}{\operatorname{mcm}}
\newcommand{\im}{\operatorname{Im}}

\title[FdM II -- Clase 05]{Algoritmo de Euclides e identidad de Bézout}
\subtitle{Cálculo efectivo del \(\mcd\)}
\author{Antonio Falcó Montesinos}
\institute{Universidad CEU Cardenal Herrera}
\date{Fundamentos de las Matemáticas II}

\begin{document}

\begin{frame}
  \titlepage
\end{frame}

\begin{frame}{Objetivos de la clase}
\begin{itemize}
  \item Calcular \(\mcd\) mediante divisiones euclídeas.
  \item Expresar el \(\mcd\) como combinación lineal entera.
  \item Aplicar la identidad de Bézout.
\end{itemize}
\end{frame}

\begin{frame}{Mapa de la clase}
\tableofcontents
\end{frame}

\section{Algoritmo de Euclides}

\begin{frame}{Algoritmo de Euclides}
\begin{block}{Idea}
\begin{itemize}
  \item Si \(a=bq+r\), entonces \(\mcd(a,b)=\mcd(b,r)\).
  \item La repetición reduce los números hasta llegar a resto cero.
  \item El último resto no nulo es el \(\mcd\).
\end{itemize}
\end{block}
\end{frame}

\begin{frame}{Algoritmo de Euclides}
\begin{block}{Ejemplo}
\begin{itemize}
  \item \(252=105\cdot 2+42\).
  \item \(105=42\cdot 2+21\).
  \item \(42=21\cdot 2+0\).
  \item Por tanto, \(\mcd(252,105)=21\).
\end{itemize}
\end{block}
\end{frame}

\begin{frame}{Algoritmo de Euclides}
\begin{block}{Ventaja}
\begin{itemize}
  \item Evita listar divisores.
  \item Funciona con enteros grandes.
  \item Prepara el cálculo de coeficientes de Bézout.
\end{itemize}
\end{block}
\end{frame}

\section{Bézout}

\begin{frame}{Bézout}
\begin{block}{Identidad}
\begin{itemize}
  \item Si \(d=\mcd(a,b)\), existen \(x,y\in\mathbb{Z}\) tales que \(d=ax+by\).
  \item La expresión \(ax+by\) se llama combinación lineal entera.
  \item Los coeficientes no son únicos.
\end{itemize}
\end{block}
\end{frame}

\begin{frame}{Bézout}
\begin{block}{Retroceso}
\begin{itemize}
  \item Se parte de las igualdades del algoritmo de Euclides.
  \item Se despeja el último resto no nulo.
  \item Se sustituyen hacia atrás hasta obtener combinación de \(a\) y \(b\).
\end{itemize}
\end{block}
\end{frame}

\section{Profundización}

\begin{frame}{Profundización}
\begin{block}{Coprimalidad}
\begin{itemize}
  \item \(a\) y \(b\) son coprimos si y solo si existen \(x,y\in\mathbb{Z}\) con \(ax+by=1\).
  \item Esta caracterización será esencial en aritmética modular.
  \item Permite detectar inversos modulares.
\end{itemize}
\end{block}
\end{frame}

\begin{frame}{Profundización}
\begin{block}{Ejercicio guiado}
\begin{itemize}
  \item Calcular \(\mcd(84,126)\).
  \item Expresarlo como \(84x+126y\).
  \item Comprobar la igualdad final sustituyendo los coeficientes.
\end{itemize}
\end{block}
\end{frame}

\begin{frame}{Profundización}
\begin{block}{Control}
\begin{itemize}
  \item ¿Cada división respeta \(0\leq r<b\)?
  \item ¿El último resto no nulo divide a los dos números iniciales?
  \item ¿La combinación lineal final da exactamente el \(\mcd\)?
\end{itemize}
\end{block}
\end{frame}


\section{Detalle y práctica}

\begin{frame}{Detalle y práctica}
\begin{block}{Algoritmo extendido}
\begin{itemize}
  \item El algoritmo de Euclides no solo calcula \(\mcd(a,b)\).
  \item Al sustituir hacia atrás, produce coeficientes \(u,v\in\mathbb{Z}\) con \(ua+vb=\mcd(a,b)\).
  \item Esta identidad es el puente hacia inversos modulares y ecuaciones diofánticas.
\end{itemize}
\end{block}
\end{frame}

\begin{frame}{Detalle y práctica}
\begin{block}{Actividad breve}
\begin{itemize}
  \item Aplicar Euclides a \(84\) y \(30\).
  \item Sustituir hacia atrás hasta obtener \(u\cdot 84+v\cdot 30=6\).
  \item Verificar que los coeficientes de Bézout no son únicos.
\end{itemize}
\end{block}
\end{frame}

\begin{frame}{Cierre}
\begin{itemize}
  \item Euclides calcula el \(\mcd\) de forma eficiente.
  \item Bézout transforma el \(\mcd\) en combinación lineal.
  \item La coprimalidad se expresa mediante una combinación igual a \(1\).
\end{itemize}
\end{frame}

\begin{frame}{Trabajo recomendado}
\begin{itemize}
  \item Releer las secciones correspondientes del manual.
  \item Identificar definiciones, símbolos nuevos y enunciados principales.
  \item Preparar dudas y ejemplos para el seminario asociado.
\end{itemize}
\end{frame}

\end{document}
